% Auto-generated by scripts/run_paper_classical_frontier_control.py
\begin{table}[t]
\centering
\footnotesize
\setlength{\tabcolsep}{4pt}
\caption{Classical-frontier control over $n=59$ measured surfaces. The table compares the default classical paper workflow, a classical least-squares unwrap, a stronger classical least-squares plus frame-normalised workflow, and a classical two-colour synthetic-wavelength baseline. The `Best classical frontier (oracle)` column is an optimistic upper bound formed by taking the per-surface minimum for each endpoint across these classical baselines and then reporting the resulting dataset median; it is therefore not a single deployable workflow. The final column shows the main hybrid branch for reference.}
\label{tab:classical_frontier_control}
\resizebox{\textwidth}{!}{%
\begin{tabular}{lcccccc}
\toprule
Endpoint & Classical default & Classical LS unwrap & Classical LSQ + norm & Classical 2-colour & Best classical frontier (oracle) & Hybrid \\
\midrule
Height RMSE (nm) & $1065.7$ & $559.2$ & $559.1$ & $1209.9$ & $559.1$ & $314.0$ \\
$|\Delta S_a|$ (nm) & $301.8$ & $366.3$ & $366.8$ & $265.4$ & $74.2$ & $129.6$ \\
$|\Delta S_q|$ (nm) & $423.4$ & $729.4$ & $728.4$ & $316.8$ & $84.7$ & $419.1$ \\
$|\Delta S_z|$ (nm) & $37163.2$ & $43202.2$ & $43198.0$ & $2278.5$ & $1368.5$ & $37457.4$ \\
\bottomrule
\end{tabular}
}%
\end{table}
